\begin{center}
\(\ast\qquad\ast\)
\end{center}

我们于是达到了可谓几何学黄金时代的时期，它大致介于蒙日《画法几何》（1795）（X）与 F. 克莱因的“埃尔朗根纲领”（1872）（XXV b）两书的出版年代之间。我们得自几何学这一突然复兴的主要进展如下：

\begin{enumerate}
\def\labelenumi{\Alph{enumi})}
\item
  无穷远元素（点、直线或平面）的概念由德萨格于十七世纪引入（VI），但在十八世纪，它几乎只是作为对术语的滥用而出现；彭赛列恢复了这一概念的地位并系统地加以使用（XII），从而把射影空间确立为一切几何现象的一般框架。
\item
  与此同时，在蒙日、尤其是彭赛列那里，完成了向复射影几何的过渡。十八世纪中偶尔使用的虚点概念，在这里（与无穷远点概念一起）得到利用，以得出独立于实仿射几何各种“图形情形”的命题。诚然，起初为支持这些革新而提出的论证仍然十分笨拙（几何学“综合”学派的拥护者们尤其如此，在那里，坐标的使用曾被视作一种污点），但人们不能不从中辨认出——在蒙日名下的“偶然关系原理”、在彭赛列名下的“连续性原理”之下——现代代数几何中“特殊化”思想的最初萌芽 (*)。
\end{enumerate}

由这些观念所得的首批成果之一，是注意到在复射影空间中，一切非退化二次曲线（相应地，二次曲面）都具有同一本性；这引导彭赛列发现了“迷向”元素：“平面上任意放置的各个圆，”他说，“因此并非如人们乍看之下可能相信的那样彼此完全独立；它们在理想意义下于无穷远处有两个公共的虚点”（（XII），第 48 页）。再往后，他又同样引入了“脐”，即所有球面公共的无穷远虚二次曲线（（XII），第 370 页）；尽管他没有特别论及球面的迷向母线，他至少明确强调了所有二次曲面都具有实的或虚的直母线这一事实（同上书，第 371 页）(*)；他的后继者们（特别是普吕克与夏尔）对这些概念的利用甚至比他本人更多，尤其是在二次曲线与二次曲面的“焦点”性质的研究中。

\begin{enumerate}
\def\labelenumi{\Alph{enumi})}
\setcounter{enumi}{2}
\item
  点变换与变换的复合这两个概念，也同样以一般的方式得到表述，并被系统地引入作为证明手段。在位移与投影之外，当时已知的变换只有寥寥几种特殊情形：拉伊尔与牛顿使用的 \(x' = a/x\)、\(y' = y/x\) 型的某些平面射影变换，来自克莱罗与欧拉的“仿射” \(x' = ax\)、\(y' = by\)，最后是牛顿、麦克劳林与布雷肯里奇那里的某些特殊的二次变换。蒙日在其《画法几何》中展示了平面投影在三维几何中所能发挥的全部作用。在彭赛列那里，一种被使用到过分程度的系统证明方法，在于借助投影把二次曲线的性质化为圆的性质（德萨格与帕斯卡已经偶尔用过这一方法）；而为了类似地从二次曲面过渡到球面，他发明了空间中射影变换的第一个例子，即“透射”（（XII），第 357 页）；最后，也正是他引入了曲线到自身的双有理变换的第一批例子。1827 年，默比乌斯（（XIII a），第 217 页）（以及 1830 年的夏尔，独立地（XV b））定义了最一般的射影线性变换；同一时期还出现了反演（参见 §10，习题 13）以及其他类型的二次变换，对它们的研究将开创双有理变换的理论，这一理论将在十九世纪下半叶得到发展。
\item
  对偶的概念充分显露出来，并被自觉地与双线性型理论联系起来。关于二次曲线的极点与极线的理论——自阿波罗尼奥斯以来仅在德萨格与拉伊尔的著作中取得过进展——被蒙日推广到二次曲面，蒙日与其学生们一起察觉到，借此手段可以把已知的定理变换为新的结果 (*)。但是，把上述种种观察确立为一般方法——即其“配极变换”理论——并把该方法变成一种特别有效的发现工具，这一功劳仍应归于彭赛列。稍后，特别是在热尔岗、普吕克、默比乌斯与夏尔那里，一般的对偶概念从它与二次型的联系（这种联系在彭赛列那里仍过于狭窄）中脱离出来。特别是，默比乌斯在考察三维空间中对偶（由一个双线性型定义）的各种可能性时，于 1833 年发现了关于交错双线性型的对偶 (XIII b) (**)，它在十九世纪尤其以“线性丛”理论的形式（参见 §10，习题 16）得到研究，并与凯莱、格拉斯曼和普吕克约在 1860 年前后引入的“直线几何”及“普吕克坐标”相联系而发展起来。
\item
  从射影几何的最初时期起，对经典几何诸性质及其与射影空间关系的深入研究，便很快导致把这些性质划分为“射影性质”与“度量性质”；把这一区分视为当时未来将成为现代结构概念的东西的最清晰表现之一，大概并不算夸张。但首先引入这一区分及这套术语的彭赛列，已经意识到联系这两类性质的东西；他在其《论著》中处理关于角的诸问题时说，角的性质“似乎不属于我们称之为射影的那类性质\ldots，然而它们却以如此简单的方式”，他说，“从构成 {[}本书{]}\ldots 之基础的诸原理导出，以致我不相信任何别的几何理论能以同时更直接、更简单的方式导出它们。如果考虑到图形的射影性质必然是图形所能具有的性质中最一般的一类，人们对这一点就丝毫不会感到惊讶；因而它们必定把一切其余的特殊性质或广延关系作为简单的推论包含在内”（（XII），第 248 页）。说实在的，在这番宣言之后，人们看到他处理角的问题时竟绕了很大的弯子——把角与二次曲线的焦点性质挂靠在一起，而不是直接让圆点发挥作用——不免有些诧异；而事实上，直到 30 年后，拉盖尔（当时还是综合理工学校的学生）才借助两条直线与同原点的两条迷向直线的交比，给出两直线间角的表达式（参见 §10，习题 5）（（XXI），卷 II，第 13 页）。最后，在凯莱 (XVIII d) 那里，一个基本思想被清楚地表达出来：平面图形的“度量”性质无非是该图形添上圆点之后的“射影”性质——这是通向“埃尔朗根纲领”的决定性里程碑。
\item
  双曲非欧几何诞生于 1830 年前后，起初多少置身于我们正概述其主线的那场运动之外。这一新几何源于本质上属于逻辑方面、触及经典几何基础的关切，其发明者们 (*) 以与欧几里得几何相同的公理化“综合”形式呈现它，且与射影几何毫无联系（在这种几何中，按照经典模式引入射影几何甚至看起来先验地就被排除了，因为唯一平行线的概念在这种几何中消失了）；大概正是由于这个原因，在很长时间里它很少引起法国、德国与英国的射影几何学派的兴趣。因此，当凯莱在上文引用的那篇基本著作 (XVIII d) 中想到把圆点（被视为一条“切向退化的”二次曲线）换成一条任意的二次曲线（他称之为“绝对形”）时，他完全没有想到把这个想法与罗巴切夫斯基–鲍耶几何联系起来，尽管他指出了他的构想如何给出两点间“距离”的新表达式，也尽管他提到了它与球面几何的联系。形势约在 1870 年前后发生变化，这时，继罗巴切夫斯基著作的传播，以及高斯著作与黎曼就职演讲的出版之后，非欧几何诸几何走到了数学现实的前沿。沿黎曼所开辟的道路，贝尔特拉米在不知道凯莱工作的情况下，于 1868 年重新发现了后者所给出的那些距离表达式，但所处的脉络完全不同：他把圆的内部视为一张常曲率曲面的像，其中测地线由直线表示 (XXIV)；两年之后，克莱因（独立于贝尔特拉米）作出了这些不同观点的综合，并以椭圆非欧空间的发现完成了这一综合 (XXV a) (**)。
\item
  在我们此处所考虑时期的后半段，出现了一个批判反思的时期，在此期间，“综合”几何的拥护者们不满足于已把坐标从证明中驱逐出去，还声称甚至在几何公理中也可以不用实数。这一学派的主要代表是冯·施陶特，他基本上完成了这一杰作 (XXIII)；此杰作在其时代乃至二十世纪相当长的时期内都备受推崇；尽管今天人们不再赋予这一层次的思想以同样的重要性——其富有成效的应用可能性已被证明相当贫乏——但仍须承认，冯·施陶特及其门徒们的努力，有助于澄清关于实或复“标量”在经典几何中所起作用的种种观念，并由此引入了在任意基体上的几何这一现代概念。
\end{enumerate}

约在 1860 年前后，“综合”几何正处于鼎盛时期，但其统治的终结已在迅速临近。在整个十八世纪始终笨重而不雅致的解析几何，在拉梅、博比利耶、柯西、普吕克与默比乌斯手中，终于获得了足以与对手平起平坐地竞争的优雅与简洁。尤其重要的是，约从 1850 年起，群与不变式的观念终于得到精确表述，逐渐占据了中心位置，并且人们认识到，经典几何的定理无非是相似变换群的不变式或协不变式之间恒等关系的表达 (*)，正如射影几何的定理表达射影群的协不变式之间的恒等式（即“合冲”）一样。这正是 F. 克莱因在其著名的“埃尔朗根纲领”(XXV b) 中以精湛技艺阐述的论点，他在其中主张放弃“综合”倾向与“解析”倾向之间徒劳无益的争论；他说，如果针对后者把特权角色给予任意坐标系的那一指责，“就以往使用坐标法时的缺陷方式而言往往确实是完全站得住脚的，那么一当涉及对这一方法的合理运用，它便不复成立\ldots{} 空间直观的领域并非解析法的禁地\ldots{}”；他还强调，“人们不可低估一个贴切的形式体系给后续研究带来的益处，因为它可以说替思考预先作了准备”（（XXV b），第 488–490 页）。

由此得到“几何学”诸定理按其源出之群的一种合理的“结构性”分类：射影几何对应线性群，度量问题对应正交群，“线性丛”的几何对应辛群。然而在这种无情的明晰之下，经典几何——代数几何与微分几何（**）除外，它们此时已确立为自治的科学——骤然枯萎，失去了全部光彩。基于使用变换的各种方法的一般化，早已使新定理的形成变得有些机械化：“今天，”夏尔于 1837 年在其 \emph{《历史概述》} 中说，“任何人都可以站出来，任取一条已知真理，把它置于各种一般变换原理之下；他将从中提取出别的真理，不同的或更一般的；而这些真理又可经受类似的运算；于是人们几乎可以无穷尽地增殖从第一条真理演绎出的新真理的数目\ldots{} 因此，在科学的现状下，谁愿意谁就可以在几何学中作出推广并有所创造；天才对于为大厦添砖加瓦已不再是不可或缺的了”（（XV b），第 268–269 页）。但随着不变量理论的进展，形势变得清楚得多：该理论终于（至少对“典型”群而言）成功地表述了一般方法，使人们原则上能以纯粹机械的方式写出一切代数协不变式及其全部“合冲”；这是一场胜利，它同时标志着经典不变量理论本身作为一个研究领域、以及“初等”几何的死亡——后者实际上已几乎沦为它的一本单纯词典。诚然，没有什么能使人们先验地预见，在可以如此随意涌现的无穷多条“定理”之中，哪些定理用适当的几何语言来表达会具有可与经典结果相比拟的简洁与优雅，而且仍存在一个有限的领域，在那里众多业余爱好者继续愉快地钻研（三角形几何、四面体几何、低次代数曲线与代数曲面的几何，等等）。但对职业数学家而言，矿藏已然枯竭，因为那里不再存在结构问题、不再存在能对数学其他部分产生反响的问题；群论与不变量理论的这一篇章，可以认为在另行通知之前业已关闭。(*)

\[
\ast \ast
\]

这样，在埃尔朗根纲领之后，欧氏几何与非欧几何从纯代数的观点来看，都成了表达双线性型理论结果的或繁或简的语言，而双线性型理论的进步则与不变量理论的进步携手并进 (*)。凡是涉及双线性型的秩的概念以及这些型与线性变换之间关系的一切内容，都由弗罗贝尼乌斯的工作最终澄清 (XXVII a)。自由 Z-模上交错型的典范表达式（§5，第 1 小节，定理 1）也应归功于弗罗贝尼乌斯 (XXVII b)；不过，斜对称行列式早在本世纪初就已出现在普法夫那里，与微分形式化为正规形相关；于 1827 年重新处理这一问题的雅可比 (XIX a)，知道奇数阶的斜对称行列式为零，并且正是他构造了普法夫式的表达式并证明它是偶数阶斜对称行列式的一个因子；但他未曾察觉后者是普法夫式的平方，这一点直到 1849 年才由凯莱建立 (XVIII b)。与二次型相伴的对称双线性型这一概念，是“极化”过程的最简单情形，而极化过程正是不变量理论的基本工具之一。以“标量积”之名，这一概念将享有巨大的成功，首先是在“向量分析”的普及者那里，然后从二十世纪起，得益于希尔伯特空间理论给它带来的意想不到的推广（见第 V 卷的历史注记）。也正是后一理论，将把算子的伴随这一概念显现出来（此前它几乎只是在线性微分方程理论中、以及在张量分析中通过协变与逆变指标在度量张量指挥棒下的舞蹈表现出来）；最后，也正是这一理论，将充分凸现埃尔米特型的概念，这一概念由埃尔米特于 1853 年在与数论研究相关的场合首先引入（（XXII），第 237 页），但直到 1925 年前后、直到复希尔伯特空间在量子理论中的应用，它一直多少处于数学主流的边缘。

另一方面，对正交群与相似变换群的研究——这两个群自十九世纪中叶起即被清楚地构想出来并被当作这样的对象处理，并且已成为二次型理论的核心——以及对其他“典型”群（线性群、辛群与酉群）的研究，扮演着越来越重要的角色。我们在此只能提一下这些群一方面在李群理论与微分几何中所起的关键作用，另一方面在二次型的算术理论中所起的关键作用（例如见 (XXXIII) 与 (XXXV)）(**)；正是由于这一情况，以及由于对偶概念向最形形色色问题的推广，才使得几乎没有哪一门现代数学理论不以这种或那种方式牵涉双线性型。无论如何我们必须指出，正是对（三维）旋转群的研究导致哈密顿发现了四元数 (XVII)；这一发现由 W. 克利福德加以推广，他于 1876 年引入了以他的名字命名的那些代数，并证明它们是四元数代数、或四元数代数与一个二次扩张的张量积 (XXVIII)。这些代数四年后被李普希茨重新发现 (XXIX)，他用它们给出了 n 个变量的正交变换的参数表示（推广了凯莱借助四元数理论对 n = 3 (XVIII a) 与 n = 4 (XVIII c) 所得到的那些表示（参见 §9，习题 15 与 16）），这些代数以及由之导出的“旋量”概念（见 (XXXII) 与 (XXXIV)），也因其在量子理论中的应用而在现代时期享有巨大的声誉。

最后，关于导致在半双线性型理论中几乎完全放弃对标量环任何限制的思想演变，我们只剩一句话要说——这是全部现代代数所共有的倾向，但在这里，它也许比在别处显现得更早。我们已经指出过在复数域上引入几何的丰硕成果（此外，在整个十九世纪，这种几何与实几何之间的混淆从未间断，有时甚至是危险的）；这里的明晰首先来自十九世纪末关于几何基础的那些公理化研究 (XXX)。在这些考察的过程中，希尔伯特及其后继者们特别是在研究各条公理之间的关系时，被引导去构造适当的反例，其中“基体”（交换或非交换）具有或多或少病态的性质，他们由此使数学家们习惯了全新类型的“几何”。从解析的观点看，伽罗瓦已经考虑过系数与变量都取值于一个有限素域的线性变换（（XVI），第 27 页）；在发展这些思想时，若尔当 (XXVI) 很自然地被引导去考虑这些体上的典型群，这些群的介入显现在数学的各个不同领域。迪克森约在 1900 年把若尔当的研究推广到一切有限域，而最近人们认识到，若尔当–迪克森理论的一大部分可以推广到绝对任意的“基体”的情形；这本质上应归功于迷向向量的一般性质与维特定理，它们在经典情形是平凡的，但直到 1936 年才对任意基体建立起来 (XXXI) (*)。

但是，在把半双线性型的研究推向越来越高的“抽象”时，把在 2 维与 3 维空间情形下来自经典几何的术语原封不动地保留下来，并把它推广到 n 维情形乃至无限维空间，被证明是极富启示性的。作为一门自主而有生命的科学已然过时的经典几何，由此被点化为当代数学的一种通用语言，其灵活与便利无可比拟。

\section*{参考文献}
\phantomsection
\addcontentsline{toc}{section}{参考文献}

\begin{enumerate}
\def\labelenumi{(\Roman{enumi})}
\item
  O. NEUGEBAUER，《古代数学科学史讲义》，第 I 卷：希腊以前的数学，柏林（Springer），1934。
\item
  J. TROPFKE，《初等数学史》，第 IV–VI 卷，柏林–莱比锡（de Gruyter），1923–24。
\item
  欧几里得《几何原本》，5 卷，J. L. Heiberg 编，莱比锡（Teubner），1883–88。
\end{enumerate}

(III bis) T. L. HEATH，《欧几里得〈几何原本〉十三卷》\ldots，3 卷，剑桥，1908。

\begin{enumerate}
\def\labelenumi{(\Roman{enumi})}
\setcounter{enumi}{3}
\item
  T. L. HEATH，《佩尔加的阿波罗尼奥斯：圆锥曲线论》，剑桥（Univ. Press），1896。
\item
  Ptolemaei Cl. Opera（托勒密全集），J. L. Heiberg 编，2 卷，莱比锡（Teubner），1898–1903。
\item
  G. DESARGUES，《全集》\ldots 第 I 卷，巴黎（Leiber），1864：Brouillon project d'une atteinte aux événements des rencontres d'un cône avec un plan，第 103-230 页。
\item
  P. FERMAT，《全集》，第 I 卷，巴黎（Gauthier-Villars），1891：a) Ad locos planos et solidos Isagoge，第 91-110 页（法文译文见同书第 III 卷，第 84-101 页）；b) Isagoge ad locos ad superficiem，第 111-117 页（法文译文见同书第 III 卷，第 102-108 页）。
\item
  L. EULER：a) Introductio in Analysin Infinitorum（Opera Omnia (1)，第 IX 卷，苏黎世–莱比锡–柏林（O. Füssli 与 B. G. Teubner），1945）；b) Theoria motus corporum solidorum seu rigidorum（Opera Omnia (2)，第 III 卷，苏黎世–莱比锡–柏林（O. Füssli 与 B. G. Teubner），1948）；c) Problema algebraicum ob affections prorsus singulares memorabile（Opera Omnia，(1)，第 VI 卷，莱比锡–柏林（Teubner），1921，第 287-315 页）；d) Formulae generales pro translatione quacunque corporum rigidorum，Novi Comm. Acad. Sc. imp. Petrop.，第 XX 卷（1776），第 189-207 页；e) De centro similitudinis.（Opera Omnia (1)，第 XXVI 卷，苏黎世（O. Füssli），1956，第 276-285 页）。
\item
  J. L. LAGRANGE，《全集》，巴黎（Gauthier-Villars），1867-1892：a) 《关于极大极小方法的研究》，第 I 卷，第 3–20 页；b) 《算术研究》，第 III 卷，第 695–795 页。
\item
  G. MONGE，《画法几何》，巴黎，1798。
\item
  C. F. GAUSS，《全集》：a) Disquisitiones arithmeticae，第 I 卷，哥廷根，1870；b) 《空间的变换》，第 VIII 卷，哥廷根，1900，第 357–362 页。
\item
  J.-V. PONCELET，《图形的射影性质论》，第 I 卷，第 2 版，巴黎（Gauthier-Villars），1865。
\item
  A. F. Möbius，《全集》，莱比锡（Hirzel），1885–87：a) 《重心演算》，第 I 卷，第 1–388 页；b) 《论空间图形之间一类特殊的对偶关系》，第 I 卷，第 489–515 页（= Crelle's Journal，第 X 卷，1833）；c) 《论解析球面几何的一种新处理》，第 II 卷，第 1–54 页。
\item
  A.-L. CAUCHY：a) 《关于无穷小演算在几何中的应用的讲义》（《全集》，(2)，第 V 卷，巴黎（Gauthier-Villars），1903）；b) 《论用以确定行星长期摄动的方程》（《全集》，(2)，第 IX 卷，巴黎（Gauthier-Villars），1891，第 174–195 页）。
\item
  M. CHASLES：a) 《关于两个物体所成系统的一般性质的注记》，Bull. de Férussac，第 XIV 卷（1830），第 321–326 页；b) 《几何方法的起源与发展的历史概述》，布鲁塞尔，1837。
\item
  E. GALOIS，《数学著作》，巴黎（Gauthier-Villars），1897。
\item
  W. R. HAMILTON，《四元数讲义》，都柏林，1853。
\item
  A. CAYLEY，《数学论文集》，剑桥，1889-1898：a) 《关于与四元数有关的若干结果》，第 I 卷，第 123-126 页（= Phil. Mag.，1845）；b) Sur les déterminants gauches，第 I 卷，第 410-413 页（= J. de Crelle，第 XXXVIII 卷（1848））；c) Recherches ultérieures sur les déterminants gauches，第 II 卷，第 202-215 页（= J. de Crelle，第 L 卷（1855））；d) 《关于代数型的第六篇论文》，第 II 卷，第 561-592 页（= Phil. Trans.，1859）。
\item
  C. G. J. JACOBI，Gesammelte Werke（《全集》），柏林（G. Reimer），1881-1891：a) 《论普法夫方法》\ldots 第 IV 卷，第 17–29 页；b) 《论一个基本的代数定理及其应用》，第 III 卷，第 593–598 页。
\item
  J. J. SYLVESTER，《数学论文集》，第 I 卷，剑桥，1904：《每个齐次二次多项式均可经实正交代换化为正负平方之和的形式这一定理的证明》，第 378–381 页（= Phil. Mag.，1852）。
\item
  E. LAGUERRE，《全集》，第 II 卷，巴黎（Gauthier-Villars），1905。
\item
  C. HERMITE，《全集》，第 I 卷，巴黎（Gauthier-Villars），1905：《论二次型理论》，第 200–263 页（= J. de Crelle，第 XLVII 卷（1854））。
\item
  K. G. V. von STÄUDT，《位置几何学的贡献》，纽伦堡，1856。
\item
  E. BELTRAMI：a) 《非欧几何解释试论》，Giorn. di Mat.，第 VI 卷（1868），第 284–312 页；b) 《常曲率空间的基本理论》，Ann. di Mat. (2)，第 II 卷（1868–69），第 232–255 页。
\item
  F. KLEIN，《数学文集》，第 I 卷，柏林（Springer），1921：a) 《论所谓非欧几何》，第 254-305 页（= Math. Ann.，第 IV 卷（1871））；b) 《关于近代几何研究的比较考察》，第 460-497 页（= Math. Ann.，第 XLIII 卷（1893））。
\item
  C. JORDAN，《置换与代数方程论》，巴黎（Gauthier-Villars），1870。
\item
  G. FROBENIUS：a) 《论线性置换与双线性型》，J. de Crelle，第 LXXXIV 卷（1878），第 1-63 页；b) 《整系数线性型的理论》，J. de Crelle，第 LXXXVI 卷（1879），第 146-208 页。
\item
  W. K. CLIFFORD，《数学论文集》，伦敦（Macmillan），1882：a) 《论几何代数的分类》，第 397-401 页；b) 《格拉斯曼广延代数的应用》，第 266-276 页（= Amer. Journ. of Math.，第 I 卷（1878））。
\item
  R. LIPSCHITZ，《关于平方和的研究》，波恩，1886。
\item
  D. HILBERT，《几何学基础》，莱比锡（Teubner），1899。
\item
  E. WITT，《任意体上的二次型理论》，J. de Crelle，第 CLXXVI 卷（1937），第 31-44 页。
\item
  E. CARTAN，《旋量理论讲义》，Actual. Sci. et Industr.，第 643 与 701 号，巴黎（Hermann），1938。
\item
  C. L. SIEGEL，《辛几何》，Amer. Journ. of Math.，第 LXV 卷（1943），第 1-86 页。
\item
  C. CHEVALLEY，《旋量的代数理论》，纽约（Columbia Univ. Press），1954。
\item
  M. EICHLER，《二次型与正交群》，柏林–哥廷根–海德堡（Springer），1952。
\end{enumerate}

\chapter*{记号索引}
\phantomsection
\addcontentsline{toc}{chapter}{记号索引}

\(b^{J}, b^{J'} : 1, 2.\)

参考编号依次表示节与目（个别情形为习题）。

\(F^{J}\)（右模 F 的标量环 B 的反自同构）：1, 2。

\(N^{0}, M^{0}\)（N、M 为子模）：1, 3。

\(u^{*}\)（u 为同态）：1, 8。

\(M(x), \mathbf{x} (x \text{ 自由模中的元素}) : 1, 10.\)

\(M(u)\)（u 为自由模到自由模的同态）：1, 10。

\(^{t}M, \dot{M}^{J} (M \text { matrix } ): 1, 10.\)

\(\mathrm{D}_{\Phi}(x_{1},\ldots ,x_{n}),\mathrm{D}_{\Phi}(\mathrm{S}):2.\)

\(\overline{\alpha}\)（\(\alpha\) 为标量）：3。

\(\operatorname{Pf}(R):5,2.\)

\(\mathbf{Sp}(\Phi), \mathbf{Sp}(2m, \mathrm{A}), \mathbf{Sp}_{2m}(\mathrm{A}) : 5, 3.\)

\(\mathbf{U}(\Phi), \mathbf{S}\mathbf{U}(\Phi) (\Phi \text{ Hermitian 形式}) : 6, 2.\)

O(Q)、SO(Q)（Q 为二次型）：6, 2。

\(\mathbf{U}(n,\mathrm{A}),\mathbf{S}\mathbf{U}(n,\mathrm{A}),\mathbf{O}(n,\mathrm{A}),\mathbf{SO}(n,\mathrm{A}):6,2.\)

\(Q^{\top}Q^{\prime}, Q \sim Q^{\prime}(Q, Q^{\prime}\text{ 二次型}): 8, 1.\)

θ(Q)（Q 为二次型）：8, 1。

\(\mathbf{T} + \mathbf{T}', a. \mathbf{T} (\mathbf{T}, \mathbf{T}'\) 为二次型的型）：8, 2。

\(\delta (\mathbb{Q})\)（Q 为二次型）：8, 2。

\(\mathbf{Q} \otimes \mathbf{Q}'\)（Q、\(\mathbf{Q}'\) 为二次型）：8, 3。

\(\mathrm{TT}^{\prime}\)（T、\(\mathrm{T}^{\prime}\) 为二次型的型）：8, 3。

\(C(Q), I(Q), \rho_{Q}, \rho, C^{+}(Q), C^{-}(Q), C^{+}, C^{-} (Q \text{ 二次型}): 9, 1.\)

\(\alpha, \beta, \mathrm{G}(f): 9, 1.\)

\(e_x, i_f, i_x^{\mathrm{F}}, \lambda_{\mathrm{F}}: 9, 2.\)

\(\lambda_{\mathbf{F}}:9,3.\)

A(Φ)（Φ 为 2 维向量空间上的对称双线性型）：10, 1。

\(\overline{u} (u\) 为正相似变换）\(i,d:10,1\)。

\((\widehat{\mathbf{D}_1,\mathbf{D}_2})\) \((\mathbf{D}_1,\mathbf{D}_2\) 为直线或半直线）：10, 3。

\[
\mathfrak {A}, \mathfrak {A} _ {0}, h, h ^ {\prime}: 1 0, 3.
\]

\(|x|, \widehat{(x,y)}(x,y\text{ 向量}):10,3.\)

cos θ、sin θ、tan θ、cot θ：10, 3。

\(\left.\right)D_{1}, D_{2}\left\{\left(D_{1}, D_{2}\right)\left(D_{1}, D_{2} \text{ 定向平面中的半直线}\right): 10, 4.\right.\)

\chapter*{术语索引}
\phantomsection
\addcontentsline{toc}{chapter}{术语索引}

参考编号依次表示小节与目（个别情形为习题）。

同态的右（左）伴随：1, 8。

半线性映射的伴随：7, 3。

二次型的克利福德代数：9, 1。

交错（矩阵）：3, 1。

两条半直线的角：10, 3 及习题 6。

仿射平面上两条射线的角：10, 3。

两条直线的角：10, 3 及习题 2。

仿射平面上两条直线的角：10, 3。

两条有向直线的角：10，习题 2。

平面上两个向量的角：10, 3。

维数 \(\geqslant 2:10,3\) 的空间中两个向量的角。

直角：10, 3 及习题 2。

旋转的角：10, 3。

正相似变换的角：10，习题 6。

平角：10, 3 及习题 6。

二次型的型环：8, 3。

维特环：8, 3。

环的反自同构：1, 2。

克利福德代数的主反自同构：9, 1。

反埃尔米特（型）：3, 1。

反埃尔米特（矩阵）：3, 1。

反对称（矩阵）：3, 1。

双线性映射：1, 1。

右（左）退化的双线性映射：1, 1。

非退化双线性映射：1, 1。

与双线性映射相伴的非退化双线性映射：1, 3。

由双线性映射经标量扩张得到的双线性映射：1, 4。

模到其克利福德代数中的典范映射：9, 1。

从埃尔米特自同态所成之集到埃尔米特型所成之集的典范映射：7, 3。

集合 \(\mathfrak{A}\) 到 \(\mathrm{S}^{+}/\mathrm{H}:10,3\) 上的典范映射。

集合 \(\mathfrak{A}\) 到 \(O^{+}:10,3\) 上的典范映射。

与双线性映射相伴的右（左）线性映射：1, 1。

与半双线性型相伴的右（左）半线性映射：1, 6。

关于一个反自同构的右（左）半双线性映射：1, 2。退化的（右、左）半双线性映射：1, 2。非退化半双线性映射：1, 2。与半双线性映射相伴的非退化半双线性映射：1, 3。由半双线性映射经标量扩张得到的半双线性映射：1, 4。与双线性映射相伴的（线性映射）：1, 1。与半双线性型相伴的（半线性映射）：1, 6。与二次型相伴的（双线性型）：3, 4。正交自同构：6, 2。n 个变量的正交自同构：6, 2。克利福德代数的主自同构：9, 1。辛自同构：5, 3。2m 个变量的辛自同构：5, 3。酉自同构：6, 2。n 个变量的酉自同构：6, 2。仿射线性丛的轴：10，习题 16。正交基：6, 1。标准正交基：6, 1。辛基：5, 1。双线性（映射）：1, 1。双线性（型）：1, 6。双模：1, 1。典范：见典范映射与典范同态。二次曲面的中心：6，习题 25 及 10，习题 12。单位圆：10，习题 2。夏尔（关系式）：10, 3。克利福德（代数）：9, 1。克利福德（群）：9, 5。线性丛：10，习题 16。条件 (C)：6, 1。条件 (C′)：6, 1。条件 (T)：4, 2。迷向锥：6，习题 23。共形（群）：10，习题 14。仿射二次曲线：6，习题 25。射影二次曲线：6，习题 23。共轭（仿射线性簇）：6，习题 25。共轭（射影线性簇）：6，习题 23。角的余弦：10, 3。角的余切：10, 3。维特分解：4, 2。右（左）退化的（双线性映射）：1, 1。右（左）退化的（半线性映射）：1, 2。退化（仿射二次曲线）：6，习题 25。退化（射影二次曲线）：6，习题 23。退化（仿射二次曲面）：6，习题 25。退化（射影二次曲面）：6，习题 23。半直线：10, 3。闭半直线：10, 3。迷向半直线：10, 3。开半直线：10, 3。位移：6, 6。二次型的维数：8, 1。正（相似变换）：6, 5。半双线性型关于一组元素的判别式：2。距离：7, 1。直（角）：10, 3。带基点直线：10，习题 2。

克利福德代数的奇元素：9, 1。迷向元素：4, 1。克利福德代数的偶元素：9, 1。奇异元素：4, 1。关于半双线性映射正交的元素：1, 3。关于二次型正交的元素：3, 4。埃尔米特自同态：7, 3。正规自同态：7, 3。ε-埃尔米特（型）：3, 1。ε-埃尔米特（矩阵）：3, 1。等价（二次型）：3, 4。等价（半双线性型）：1, 6。二次型的定义空间：8, 1。欧几里得空间：6, 6。埃尔米特空间：6, 6。定向空间：10, 3。欧几里得（空间）：6, 6。半双线性型到张量幂的扩张：1, 9。半双线性型到外幂的扩张：1, 9。二次型的外（直和）：3, 4。二次模的外（直和）：3, 4。弱正交（子空间）：3，习题 11。闭（角域）：10, 4。闭（半直线）：10, 3。余弦函数：10, 3。余切函数：10, 3。正弦函数：10, 3。正切函数：10, 3。三角函数：10, 3。反埃尔米特型：3, 1。双线性型：1, 6。与二次型相伴的双线性型：3, 4。埃尔米特双线性型：3, 1。埃尔米特型：3, 1。负（正）埃尔米特型：7, 1。双线性（半双线性）型的逆型：1, 7。度量型：6, 6。中性型：4, 2 及 8, 1。二次型：3, 4。负（正）二次型：7, 1。非退化二次型：3, 4。由二次型经标量扩张得到的二次型：3, 4。半双线性型：1, 6。等价的二次型：3, 4。等价的半双线性型：1, 6。格拉姆–施密特（正交化方法）：6, 1。克利福德群：9, 5。约化克利福德群：9, 5。特殊克利福德群：9, 5。旋转群：9, 5。二次型的型群：8, 2。维特群：8, 2。与 Q 相伴的正交群：6, 2。n 个变量的正交群：6, 2。特殊正交群：6, 2。特殊酉群：6, 2。与 Φ 相伴的辛群：5, 3。2m 个变量的辛群：5, 3。与 Φ 相伴的酉群：6, 2。

\begin{Shaded}
\begin{Highlighting}[]
\NormalTok{n 个变量的酉群 : 6, 2.}
\NormalTok{埃尔米特（自同态） : 7, 3.}
\NormalTok{埃尔米特（空间） : 6, 6.}
\NormalTok{埃尔米特（型） : 3, 1.}
\NormalTok{埃尔米特（矩阵） : 3, 1.}
\NormalTok{\textquotesingle{}典范同态\textquotesingle{}——从 E 的子{-}模的克利福德代数到 E 的克利福德代数中 : 9, 1.\textquotesingle{}}
\NormalTok{度量同态 : 4, 3.}
\NormalTok{两球面的根超平面 : 10, 习题 12.}
\NormalTok{双线性映射的原像 : 1, 1.\textquotesingle{}}
\NormalTok{半双线性映射的原像 : 1, 2.\textquotesingle{}}
\NormalTok{克利福德代数中的奇（元素） : 9, 1.}
\NormalTok{二次型（半双线性型）的指数 : 4, 2.\textquotesingle{}}
\NormalTok{迪克森不变量 : 9, 习题 9.}
\NormalTok{逆（型） : 1, 7.}
\NormalTok{逆（相似变换） : 6, 5.}
\NormalTok{反演 : 10, 习题 13.}
\NormalTok{球面 C 的反演 : 10, 习题 13.}
\NormalTok{线性群中的对合 : 6, 3.}
\NormalTok{迷向（半{-}直线） : 10, 3.}
\NormalTok{迷向（元素） : 4, 1.}
\NormalTok{迷向（子{-}模） : 4, 1.}
\NormalTok{迷向（线性簇） : 6, 6.}
\NormalTok{拉盖尔（公式） : 10, 习题 5.}
\NormalTok{惯性定律 : 7, 2.\textquotesingle{}}
\NormalTok{向量的长度 : 10, 3.\textquotesingle{}}
\NormalTok{交错矩阵 : 3, 1.}
\NormalTok{反埃尔米特矩阵 : 3, 1.}
\NormalTok{反对称矩阵 : 3, 1.}
\NormalTok{满足 (G)、(D)、(GD) 或 (DG) 的映射的矩阵 : 1, 10.\textquotesingle{}}
\NormalTok{双线性映射的矩阵 : 1, 1.\textquotesingle{}}
\NormalTok{半双线性映射的矩阵 : 1, 2.\textquotesingle{}}
\NormalTok{ε{-}埃尔米特矩阵 : 3, 1.}
\NormalTok{埃尔米特矩阵 : 3, 1.}
\NormalTok{正规矩阵 : 7, 习题 17.}
\NormalTok{正交矩阵 : 6, 2.}
\NormalTok{对称矩阵 : 3, 1.}
\NormalTok{辛矩阵 : 5, 3.}
\NormalTok{酉矩阵 : 6, 2.}
\NormalTok{度量（型） : 6, 6.}
\NormalTok{度量（同态） : 4, 3.}
\NormalTok{二次模 : 3, 4.}
\NormalTok{相似变换的乘子 : 6, 5 及 6.\textquotesingle{}}
\NormalTok{负（n{-}向量） : 10, 3.}
\NormalTok{负（埃尔米特型） : 7, 1.}
\NormalTok{负（二次型） : 7, 1.}
\NormalTok{中性（二次型） : 8, 1.}
\NormalTok{中性（半双线性型） : 4, 2.}
\NormalTok{非退化（双线性映射） : 1, 1.}
\NormalTok{非退化（半双线性映射） : 1, 2.}
\NormalTok{非退化（二次型） : 3, 4.}
\NormalTok{正规（自同态） : 7, 3.}
\NormalTok{旋量范数 : 9, 5.}
\NormalTok{二次模的核 : 3, 4.\textquotesingle{}}
\NormalTok{负（正）n{-}向量 : 10, 3.}
\NormalTok{空间的定向 : 10, 3.\textquotesingle{}}
\NormalTok{定向（空间） : 10, 3.}
\NormalTok{正交（自同构） : 6, 2.}
\NormalTok{正交（群） : 6, 2.}
\NormalTok{正交（特殊群） : 6, 2.}
\end{Highlighting}
\end{Shaded}

正交（投影算子）：6, 3。与子模正交的（子模）：1, 3。与线性簇正交的（向量）：6, 6。正交（基）：6, 1。正交（矩阵）：6, 2。到线性簇上的正交（投影）：6, 6。正交（变换）：6, 2。与线性簇正交的（簇）：6, 6。正交（部分）：1, 3 及 3, 4。正交（线性簇）：6, 6。格拉姆–施密特正交化（方法）：6, 1。正交（元素）：1, 3 及 3, 4。标准正交（基）：6, 1。开（角域）：10, 4。开（半直线）：10, 3。克利福德代数中的偶（元素）：9, 1。正交部分：1, 3 及 3, 4。垂直（线性簇）：6，习题 22。交错矩阵的普法夫式：5, 2。平（角）：10, 3。球极投影的视点：10，习题 14。线性簇关于二次曲面的极簇：6，习题 23 及 25。超平面关于二次曲面的极点：6，习题 23 及 25。反演的极点：10，习题 13。正（n-向量）：10, 3。正（埃尔米特型）：7, 1。正（二次型）：7, 1。主（反自同构）：9, 1。主（自同构）：9, 1。二次型的张量积：8, 3。半双线性型的张量积：1, 9。正交投影算子：6, 3。到线性簇上的正交投影：6, 6。球极投影：10，习题 14。伪判别式：9，习题 9。点关于球面的幂：10，习题 12。反演的幂：10，习题 13。勾股（定理）：6, 6。二次（型）：3, 4。二次（模）：3, 4。仿射二次曲面：6，习题 25。射影二次曲面：6，习题 23。双线性（半双线性）型的秩：1, 6。球面的半径：10，习题 12。约化（克利福德群）：9, 5。夏尔关系式：10, 3。旋量表示：9, 4。旋转：9, 5。角为 φ 的旋转：10, 3 及习题 2。旋转（群）：9, 5。闭角域：10, 4。开角域：10, 4。平角域：10，习题 8。凹角域：10，习题 8。凸角域：10，习题 8。半旋量：9, 4。半双线性（映射）：1, 2。半双线性（型）：1, 6。埃尔米特型的符号差：7, 2。相似变换：6, 5 及 6。

正相似变换：6, 5 及 9，习题 9。

逆相似变换：6, 5。

奇异（元素）：4, 1。

奇异（子模）：4, 1。

角的正弦：10, 3。

双线性（半双线性）映射的直和：1, 3。

二次型（二次模）的外直和：3, 4。

迷向子模：4, 1。

与子模正交的子模：1, 3。

奇异子模：4, 1。

全迷向子模：4, 1。

全奇异子模：4, 1。

球面：10，习题 12。

正交球面：10，习题 12。

旋量：9, 4。

半直线的正向序列：10, 4。

关于超平面的对称（变换）：6, 4。

关于向量子空间的对称（变换）：6, 3。

关于仿射线性簇的对称（变换）：6, 6。

对称（矩阵）：3, 1。

辛（自同构）：5, 3。

辛（基）：5, 1。

辛（群）：5, 3。

辛（矩阵）：5, 3。

辛（变换）：5, 3。

角的正切：10, 3。

与二次曲面相切的（线性簇）：6，习题 23 及 25。

勾股定理：6, 6。

维特定理：4, 3。

全迷向（子模）：4, 1。

与子模完全正交的（子模）：1, 3。

与线性簇完全正交的（簇）：6, 6。

全奇异（子模）：4, 1。

正交变换：6, 2。

辛变换：5, 3。

酉变换：6, 2。

三角（函数）：10, 3。

二次型的型：8, 1。

酉（自同构）：6, 2。

酉（群）：6, 2。

酉（特殊群）：6, 2。

酉（矩阵）：6, 2。

酉（变换）：6, 2。

迷向线性簇：6, 6。

与线性簇正交的线性簇：6, 6。

全迷向线性簇：6, 6。

正交的线性簇：6, 6。

与线性簇正交的向量：6, 6。

维特（环）：8, 3。

维特（分解）：4, 2。

维特（群）：8, 2。

维特（定理）：4, 3。

\chapter*{\texorpdfstring{第\chinese{chapter}章定义汇编}{第九章定义汇编}}
\phantomsection
\addcontentsline{toc}{chapter}{\texorpdfstring{第\chinese{chapter}章定义汇编}{第九章定义汇编}}

\section*{半双线性型：}

设 A 是一个环， \(\xi \to \overline{\xi}\) 是 A 的一个对合反自同构，即 A 到其自身的一个双射，使得 \(\overline{(\xi + \eta)} = \overline{\xi} + \overline{\eta}\) 、 \(\overline{(\xi\eta)} = \overline{\eta} . \overline{\xi}\) 、 \(\overline{\overline{\xi}} = \xi\) 对 A 中一切 \(\xi, \eta\) 成立。左 A-模 E 上的半双线性型 \(\Phi\) 是指 E × E 到 A 中的一个映射，使得

\[
\begin{array}{c c} \Phi (x + x ^ {\prime}, y) = \Phi (x, y) + \Phi (x ^ {\prime}, y), & \Phi (x, y + y ^ {\prime}) = \Phi (x, y) + \Phi (x, y ^ {\prime}) \\ \Phi (\alpha x, y) = \alpha \Phi (x, y), & \Phi (x, \alpha y) = \Phi (x, y) \bar {\alpha} \end{array}
\]

对 E 中一切 \(\alpha\in\mathbf{A},x,x^{\prime},y,y^{\prime}\) 成立。

设 \(\varepsilon\) 为 A 的中心的一个元素。称 E 上的半双线性型 \(\Phi\) 为 \(\varepsilon\)-埃尔米特型，如果 \(\Phi(y, x) = \varepsilon \overline{\Phi(x, y)}\) 对一切 \(x \in E\) 、 \(y \in E\) 成立；若 \(\varepsilon = 1\) （相应地 \(\varepsilon = -1\) ），则称 \(\Phi\) 为埃尔米特型（相应地，反埃尔米特型）。

当反自同构 \(\xi\to\overline{\xi}\) 为恒等映射（这蕴含环 A 是交换环）时，相应的半双线性型就是双线性型。这时称“对称”以代替“埃尔米特”，称“反对称”以代替“反埃尔米特”（参见第 III 章）。双线性型 \(\Phi\) 若满足 \(\Phi(x,x)=0\) ，则称为交错型；此时它是反对称的，而当 A 是特征 \(\neq2\) 的域时，其逆亦真（参见第 III 章）。

称 \(\varepsilon\)-埃尔米特型满足条件 (T)，如果对一切 \(x \in E\) ，存在 \(\lambda \in A\) 使 \(\Phi(x, x) = \lambda + \varepsilon \overline{\lambda}\) 。当 \(\varepsilon = 1\) 且 \(A\) 是特征 \(\neq 2\) 的域，或当 \(\Phi\) 是交错型时，这一条件总得到满足。

\section*{二次型：}

设 A 是一个交换环。A-模 E 上的二次型 Q 是指 E 到 A 中的一个映射，使得 \(\mathrm{Q}(\alpha x) = \alpha^{2}\mathrm{Q}(x)\) 对 \(\alpha \in A, x \in E\) 成立，并且使映射

\[
(x, y) \rightarrow \Phi (x, y) = Q (x + y) - Q (x) - Q (y)
\]

是一个双线性型（必然是对称的），称为与二次型 Q 相伴的双线性型。我们有 \(\Phi(x, x) = 2Q(x)\) ；反之，对 E 上的每个双线性型 \(\Psi\) ， \(x \to \Psi(x, x)\) 是 E 上的一个二次型；

因而当 A 是特征 ≠ 2 的域时，E 上的二次型与对称双线性型彼此双射地对应。

\section*{正交元素：}

设 \(\Phi\) 是左 A-模 E 上的一个 \(\varepsilon\)-埃尔米特型。称 E 的两个元素 \(x, y\) （关于 \(\Phi\) ）正交，如果 \(\Phi(x, y) = 0\) ；这一关系关于 \(x\) 与 \(y\) 是对称的。对 E 的每个子模 M，与 M 的所有元素正交的那些 \(x \in E\) 所成之集是一个子模，记作 \(\mathbf{M}^{\circ}\) ，称为子模 M 的正交补。称 \(\Phi\) 为非退化的，如果 \(\mathbf{E}^{\circ} = \{0\}\) 。当 E 是有限维向量空间且 \(\Phi\) 非退化时，对 E 的每个子空间 M，有 codim \(\mathbf{M}^{\circ} = \dim \mathbf{M}\) 及 \(\mathbf{M}^{\circ\circ} = \mathbf{M}\) ；此外，对 E 的每一对子空间 M、N，有

\[
(\mathbf {M} + \mathbf {N}) ^ {0} = \mathbf {M} ^ {0} \cap \mathbf {N} ^ {0}, \quad (\mathbf {M} \cap \mathbf {N}) ^ {0} = \mathbf {M} ^ {0} + \mathbf {N} ^ {0}.
\]

设环 A 是交换环，Q 是 A-模 E 上的一个二次型，\(\Phi\) 是与 Q 相伴的双线性型。称 E 的两个元素（关于 Q）正交，如果它们关于 \(\Phi\) 正交；子模 M（关于 Q）的正交补就是正交补 M \(^{o}\) ，即 M 关于 \(\Phi\) 的正交补。若 \(\Phi\) 非退化，则称 Q 非退化。

\section*{迷向元素与奇异元素：}

设 \(\Phi\) 是左 A-模 E 上的一个 \(\varepsilon\)-埃尔米特型。称元素 \(x \in E\) 为迷向的，如果 \(\Phi(x, x) = 0\) ；称 E 的子模 M 为迷向的，如果 \(\mathbf{M} \cap \mathbf{M}^{\circ} \neq \{0\}\) （换言之，如果 \(\Phi\) 在 M 上的限制是退化的）；称 M 为全迷向的，如果 \(\mathbf{M} \subset \mathbf{M}^{\circ}\) （换言之，如果 \(\Phi\) 在 M 上的限制为零）。对 E 的每个子模 M，\(\mathbf{M} \cap \mathbf{M}^{\circ}\) 都是全迷向的。若 E 是有限维向量空间且 \(\Phi\) 非退化，则对 E 的子空间 M，下列三个条件等价： \(1^{\circ}\) M 非迷向； \(2^{\circ}\) M \(^{\circ}\) 非迷向； \(3^{\circ}\) E 是 M 与 M \(^{\circ}\) 的直和。设环 A 是交换环，Q 是 A-模 E 上的一个二次型，\(\Phi\) 是与 Q 相伴的双线性型。称元素 \(x \in E\) 为奇异的，如果 \(\mathrm{Q}(x) = 0\) ；称 E 的子模 M 为奇异的（相应地，全奇异的），如果 \(\mathbf{M} \cap \mathbf{M}^{\circ}\) 含有一个 \(\neq 0\) 的奇异元素（相应地，如果 Q 在 M 上的限制为零）。每个奇异元素（相应地，奇异子模、全奇异子模）都是迷向的（相应地，迷向的、全迷向的）；当 A 是特征 \(\neq 2\) 的域时，其逆亦真。
